\chapter*{\centering РЕФЕРАТ}
\addcontentsline{toc}{chapter}{РЕФЕРАТ}

Настоящая работа посвящена воспроизведению и анализу методов
 нейросетевого моделирования свойств никелевых жаропрочных сплавов. 
 Рассмотрены две задачи: прогнозирование температуры сольвуса
  $\gamma'$-фазы по химическому составу сплава и моделирование
 зависимости прочности никелевых суперсплавов от параметра 
 Ларсона--Миллера.

В первой части работы исследуется подход к прогнозированию
 температуры сольвуса с использованием искусственной нейронной 
 сети с дифференциальным слоем. Особое внимание уделено подготовке
  входных данных: учету априорной информации о физической роли 
  элементов в структуре жаропрочных сплавов. Такой подход 
  позволяет представить состав сплава в более физически 
  обоснованной форме и повысить интерпретируемость модели.

Во второй части рассматривается модель прогнозирования жаропрочности 
никелевых суперсплавов на основе байесовской регуляризованной 
искусственной нейронной сети BRANN. В качестве входных признаков 
используются химический состав и нормированный параметр 
Ларсона--Миллера, объединяющий температуру и время испытания. 
Для повышения устойчивости обучения прочность предварительно 
преобразуется в логарифмическую форму.

В работе описаны теоретические основы, подготовка данных, 
структура применяемых нейронных сетей, особенности реализации 
вычислительных экспериментов и основные результаты моделирования. 
Показано, что использование методов машинного обучения позволяет 
учитывать сложные нелинейные зависимости между химическим составом, 
температурно-временными условиями и свойствами никелевых жаропрочных 
сплавов.

\textbf{Ключевые слова:} никелевые жаропрочные сплавы, суперсплавы, 
температура сольвуса, $\gamma'$-фаза, искусственная нейронная сеть, 
дифференциальный слой, BRANN, байесовская регуляризация, параметр 
Ларсона--Миллера, жаропрочность.
